%! Author = lili
%! Date = 4/27/26


\section{Vorlesung 7}
Die Verteilung charakterisiert eine \Gls{zv} $X$ eindeutig. \\
Sind $X$ und $Y$  zwei \Glspl{zv} mit mit derselben Verteilung, so gilt
\[
    \pe(X \in A) = \pe_X(A) = \pe_Y(A) = \pe(Y \in A)
\]
für jedes $A$. \\
Es ist nicht immer klar, welche Verteilung zu einem sinnvollen Modell führt.
Auch sind die Parameter der Verteilung nicht unbedingt bekannt. \\
\texttt{BSP Verteilungen Folie 3} \\
Wie können wir trotzdem Vorhersagen treffen? \\ \newline

\paragraph{Typische Fragestellungen:} Wie \dots \textit{im Mittel}?

\begin{bsp}
    \textbf{1 $\times$ Würfeln mit einem fairem Würfel} \\
    Da herrscht eine Gleichverteilung, daher ist die erwartete Augenzahl ist das arithmetische Mittel. \\
    \texttt{Ganzen Bsp. Folie 5}
\end{bsp}

\begin{bsp}
    \textbf{Spiel, bei dem nicht alle Ausgänge gleichwahrscheinlich sind} \\
    1 $\times$ Würfeln eines fairen Würfels \\
    Auszahlung:
    \begin{itemize}
        \item[] $k$ Euro, falls 2 gewürfelt
        \item[] 6 Euro, falls $2, \cdots , 6$ gewürfelt
    \end{itemize}
    \glssymbol{eraum} = $\kl{1, \dots, 6}$ mit Gleichverteilung \pe  \\
    Die Auszahlung wird beschrieben durch die \Gls{zv} $Y$ mit
    \[
        Y(\glssymbol{elementarereignis}) = k \cdot \chi_{\kl{1}}(\ele) + 6 \cdot \chi_{\kl{2, \dots, 6}}(\ele) \quad \faoio
    \]
    Erwartete Auszahlung als arithmetisches Mittel:
    \mab{\E{Y} = \sum_{\oio} Y(\ele) \glssymbol{wfk}(\ele) = \frac{k+ 5 \cdot 6}{6} = 6 \cdot \frac{5}{6} + k \cdot \frac{1}{6}  = 6 \cdot \pe(Y=6) + k \cdot \pe(Y=k) = \sum_{I\in Y(\glssymbol{eraum})} I \cdot \pe(Y=I)}
\end{bsp}

\begin{defi}
    Der \Gls{ew} einer \zv $X$ mit abzählbarem Bildbereich $X(\eraum) \in \mathbb{R}$ ist gegeben durch
    \mab{\E{X} = \mathbb{E} X = \sum_{x\in X(\eraum)} x \cdot \pe (X=x) }
    insofern die Reihe \glsverb{abskonv}, also \ma{ \sum_{x\in X(\eraum)} |x| \cdot \pe (X=x) < \infty}. \\
    Wir sprechen in diesem Fall von der Existenz des \Gls{ew}s. \\ \newline
    Der \Gls{ew} einer \zv $X$ mit einer \dichte \glssymbol{dichte} ist gegeben durch
    \mab{\E X = \int_{- \infty}^{\infty} x \glssymbol{dichte}(x) \dx = \lim_{N \rightarrow \infty} \int_{-N}^{N} x \glssymbol{dichte}(x) \dx }
    insofern \ma{\int_{- \infty}^{\infty} |x| \glssymbol{dichte}(x) \dx < \infty}. \\
    Wir sprechen in diesem Fall von der Existenz des Erwartungswerts. \\
    \texttt{Hinweis Folie 7}
\end{defi}
